\documentclass[11pt]{article}

    \usepackage[breakable]{tcolorbox}
    \usepackage{parskip} % Stop auto-indenting (to mimic markdown behaviour)
    

    % Basic figure setup, for now with no caption control since it's done
    % automatically by Pandoc (which extracts ![](path) syntax from Markdown).
    \usepackage{graphicx}
    % Keep aspect ratio if custom image width or height is specified
    \setkeys{Gin}{keepaspectratio}
    % Maintain compatibility with old templates. Remove in nbconvert 6.0
    \let\Oldincludegraphics\includegraphics
    % Ensure that by default, figures have no caption (until we provide a
    % proper Figure object with a Caption API and a way to capture that
    % in the conversion process - todo).
    \usepackage{caption}
    \DeclareCaptionFormat{nocaption}{}
    \captionsetup{format=nocaption,aboveskip=0pt,belowskip=0pt}

    \usepackage{float}
    \floatplacement{figure}{H} % forces figures to be placed at the correct location
    \usepackage{xcolor} % Allow colors to be defined
    \usepackage{enumerate} % Needed for markdown enumerations to work
    \usepackage{geometry} % Used to adjust the document margins
    \usepackage{amsmath} % Equations
    \usepackage{amssymb} % Equations
    \usepackage{textcomp} % defines textquotesingle
    % Hack from http://tex.stackexchange.com/a/47451/13684:
    \AtBeginDocument{%
        \def\PYZsq{\textquotesingle}% Upright quotes in Pygmentized code
    }
    \usepackage{upquote} % Upright quotes for verbatim code
    \usepackage{eurosym} % defines \euro

    \usepackage{iftex}
    \ifPDFTeX
        \usepackage[T1]{fontenc}
        \IfFileExists{alphabeta.sty}{
              \usepackage{alphabeta}
          }{
              \usepackage[mathletters]{ucs}
              \usepackage[utf8x]{inputenc}
          }
    \else
        \usepackage{fontspec}
        \usepackage{unicode-math}
    \fi

    \usepackage{fancyvrb} % verbatim replacement that allows latex
    \usepackage{grffile} % extends the file name processing of package graphics
                         % to support a larger range
    \makeatletter % fix for old versions of grffile with XeLaTeX
    \@ifpackagelater{grffile}{2019/11/01}
    {
      % Do nothing on new versions
    }
    {
      \def\Gread@@xetex#1{%
        \IfFileExists{"\Gin@base".bb}%
        {\Gread@eps{\Gin@base.bb}}%
        {\Gread@@xetex@aux#1}%
      }
    }
    \makeatother
    \usepackage[Export]{adjustbox} % Used to constrain images to a maximum size
    \adjustboxset{max size={0.9\linewidth}{0.9\paperheight}}

    % The hyperref package gives us a pdf with properly built
    % internal navigation ('pdf bookmarks' for the table of contents,
    % internal cross-reference links, web links for URLs, etc.)
    \usepackage{hyperref}
    % The default LaTeX title has an obnoxious amount of whitespace. By default,
    % titling removes some of it. It also provides customization options.
    \usepackage{titling}
    \usepackage{longtable} % longtable support required by pandoc >1.10
    \usepackage{booktabs}  % table support for pandoc > 1.12.2
    \usepackage{array}     % table support for pandoc >= 2.11.3
    \usepackage{calc}      % table minipage width calculation for pandoc >= 2.11.1
    \usepackage[inline]{enumitem} % IRkernel/repr support (it uses the enumerate* environment)
    \usepackage[normalem]{ulem} % ulem is needed to support strikethroughs (\sout)
                                % normalem makes italics be italics, not underlines
    \usepackage{soul}      % strikethrough (\st) support for pandoc >= 3.0.0
    \usepackage{mathrsfs}
    

    
    % Colors for the hyperref package
    \definecolor{urlcolor}{rgb}{0,.145,.698}
    \definecolor{linkcolor}{rgb}{.71,0.21,0.01}
    \definecolor{citecolor}{rgb}{.12,.54,.11}

    % ANSI colors
    \definecolor{ansi-black}{HTML}{3E424D}
    \definecolor{ansi-black-intense}{HTML}{282C36}
    \definecolor{ansi-red}{HTML}{E75C58}
    \definecolor{ansi-red-intense}{HTML}{B22B31}
    \definecolor{ansi-green}{HTML}{00A250}
    \definecolor{ansi-green-intense}{HTML}{007427}
    \definecolor{ansi-yellow}{HTML}{DDB62B}
    \definecolor{ansi-yellow-intense}{HTML}{B27D12}
    \definecolor{ansi-blue}{HTML}{208FFB}
    \definecolor{ansi-blue-intense}{HTML}{0065CA}
    \definecolor{ansi-magenta}{HTML}{D160C4}
    \definecolor{ansi-magenta-intense}{HTML}{A03196}
    \definecolor{ansi-cyan}{HTML}{60C6C8}
    \definecolor{ansi-cyan-intense}{HTML}{258F8F}
    \definecolor{ansi-white}{HTML}{C5C1B4}
    \definecolor{ansi-white-intense}{HTML}{A1A6B2}
    \definecolor{ansi-default-inverse-fg}{HTML}{FFFFFF}
    \definecolor{ansi-default-inverse-bg}{HTML}{000000}

    % common color for the border for error outputs.
    \definecolor{outerrorbackground}{HTML}{FFDFDF}

    % commands and environments needed by pandoc snippets
    % extracted from the output of `pandoc -s`
    \providecommand{\tightlist}{%
      \setlength{\itemsep}{0pt}\setlength{\parskip}{0pt}}
    \DefineVerbatimEnvironment{Highlighting}{Verbatim}{commandchars=\\\{\}}
    % Add ',fontsize=\small' for more characters per line
    \newenvironment{Shaded}{}{}
    \newcommand{\KeywordTok}[1]{\textcolor[rgb]{0.00,0.44,0.13}{\textbf{{#1}}}}
    \newcommand{\DataTypeTok}[1]{\textcolor[rgb]{0.56,0.13,0.00}{{#1}}}
    \newcommand{\DecValTok}[1]{\textcolor[rgb]{0.25,0.63,0.44}{{#1}}}
    \newcommand{\BaseNTok}[1]{\textcolor[rgb]{0.25,0.63,0.44}{{#1}}}
    \newcommand{\FloatTok}[1]{\textcolor[rgb]{0.25,0.63,0.44}{{#1}}}
    \newcommand{\CharTok}[1]{\textcolor[rgb]{0.25,0.44,0.63}{{#1}}}
    \newcommand{\StringTok}[1]{\textcolor[rgb]{0.25,0.44,0.63}{{#1}}}
    \newcommand{\CommentTok}[1]{\textcolor[rgb]{0.38,0.63,0.69}{\textit{{#1}}}}
    \newcommand{\OtherTok}[1]{\textcolor[rgb]{0.00,0.44,0.13}{{#1}}}
    \newcommand{\AlertTok}[1]{\textcolor[rgb]{1.00,0.00,0.00}{\textbf{{#1}}}}
    \newcommand{\FunctionTok}[1]{\textcolor[rgb]{0.02,0.16,0.49}{{#1}}}
    \newcommand{\RegionMarkerTok}[1]{{#1}}
    \newcommand{\ErrorTok}[1]{\textcolor[rgb]{1.00,0.00,0.00}{\textbf{{#1}}}}
    \newcommand{\NormalTok}[1]{{#1}}

    % Additional commands for more recent versions of Pandoc
    \newcommand{\ConstantTok}[1]{\textcolor[rgb]{0.53,0.00,0.00}{{#1}}}
    \newcommand{\SpecialCharTok}[1]{\textcolor[rgb]{0.25,0.44,0.63}{{#1}}}
    \newcommand{\VerbatimStringTok}[1]{\textcolor[rgb]{0.25,0.44,0.63}{{#1}}}
    \newcommand{\SpecialStringTok}[1]{\textcolor[rgb]{0.73,0.40,0.53}{{#1}}}
    \newcommand{\ImportTok}[1]{{#1}}
    \newcommand{\DocumentationTok}[1]{\textcolor[rgb]{0.73,0.13,0.13}{\textit{{#1}}}}
    \newcommand{\AnnotationTok}[1]{\textcolor[rgb]{0.38,0.63,0.69}{\textbf{\textit{{#1}}}}}
    \newcommand{\CommentVarTok}[1]{\textcolor[rgb]{0.38,0.63,0.69}{\textbf{\textit{{#1}}}}}
    \newcommand{\VariableTok}[1]{\textcolor[rgb]{0.10,0.09,0.49}{{#1}}}
    \newcommand{\ControlFlowTok}[1]{\textcolor[rgb]{0.00,0.44,0.13}{\textbf{{#1}}}}
    \newcommand{\OperatorTok}[1]{\textcolor[rgb]{0.40,0.40,0.40}{{#1}}}
    \newcommand{\BuiltInTok}[1]{{#1}}
    \newcommand{\ExtensionTok}[1]{{#1}}
    \newcommand{\PreprocessorTok}[1]{\textcolor[rgb]{0.74,0.48,0.00}{{#1}}}
    \newcommand{\AttributeTok}[1]{\textcolor[rgb]{0.49,0.56,0.16}{{#1}}}
    \newcommand{\InformationTok}[1]{\textcolor[rgb]{0.38,0.63,0.69}{\textbf{\textit{{#1}}}}}
    \newcommand{\WarningTok}[1]{\textcolor[rgb]{0.38,0.63,0.69}{\textbf{\textit{{#1}}}}}


    % Define a nice break command that doesn't care if a line doesn't already
    % exist.
    \def\br{\hspace*{\fill} \\* }
    % Math Jax compatibility definitions
    \def\gt{>}
    \def\lt{<}
    \let\Oldtex\TeX
    \let\Oldlatex\LaTeX
    \renewcommand{\TeX}{\textrm{\Oldtex}}
    \renewcommand{\LaTeX}{\textrm{\Oldlatex}}
    % Document parameters
    % Document title
    \title{fundamentals}
    
    
    
    
    
    
    
% Pygments definitions
\makeatletter
\def\PY@reset{\let\PY@it=\relax \let\PY@bf=\relax%
    \let\PY@ul=\relax \let\PY@tc=\relax%
    \let\PY@bc=\relax \let\PY@ff=\relax}
\def\PY@tok#1{\csname PY@tok@#1\endcsname}
\def\PY@toks#1+{\ifx\relax#1\empty\else%
    \PY@tok{#1}\expandafter\PY@toks\fi}
\def\PY@do#1{\PY@bc{\PY@tc{\PY@ul{%
    \PY@it{\PY@bf{\PY@ff{#1}}}}}}}
\def\PY#1#2{\PY@reset\PY@toks#1+\relax+\PY@do{#2}}

\@namedef{PY@tok@w}{\def\PY@tc##1{\textcolor[rgb]{0.73,0.73,0.73}{##1}}}
\@namedef{PY@tok@c}{\let\PY@it=\textit\def\PY@tc##1{\textcolor[rgb]{0.24,0.48,0.48}{##1}}}
\@namedef{PY@tok@cp}{\def\PY@tc##1{\textcolor[rgb]{0.61,0.40,0.00}{##1}}}
\@namedef{PY@tok@k}{\let\PY@bf=\textbf\def\PY@tc##1{\textcolor[rgb]{0.00,0.50,0.00}{##1}}}
\@namedef{PY@tok@kp}{\def\PY@tc##1{\textcolor[rgb]{0.00,0.50,0.00}{##1}}}
\@namedef{PY@tok@kt}{\def\PY@tc##1{\textcolor[rgb]{0.69,0.00,0.25}{##1}}}
\@namedef{PY@tok@o}{\def\PY@tc##1{\textcolor[rgb]{0.40,0.40,0.40}{##1}}}
\@namedef{PY@tok@ow}{\let\PY@bf=\textbf\def\PY@tc##1{\textcolor[rgb]{0.67,0.13,1.00}{##1}}}
\@namedef{PY@tok@nb}{\def\PY@tc##1{\textcolor[rgb]{0.00,0.50,0.00}{##1}}}
\@namedef{PY@tok@nf}{\def\PY@tc##1{\textcolor[rgb]{0.00,0.00,1.00}{##1}}}
\@namedef{PY@tok@nc}{\let\PY@bf=\textbf\def\PY@tc##1{\textcolor[rgb]{0.00,0.00,1.00}{##1}}}
\@namedef{PY@tok@nn}{\let\PY@bf=\textbf\def\PY@tc##1{\textcolor[rgb]{0.00,0.00,1.00}{##1}}}
\@namedef{PY@tok@ne}{\let\PY@bf=\textbf\def\PY@tc##1{\textcolor[rgb]{0.80,0.25,0.22}{##1}}}
\@namedef{PY@tok@nv}{\def\PY@tc##1{\textcolor[rgb]{0.10,0.09,0.49}{##1}}}
\@namedef{PY@tok@no}{\def\PY@tc##1{\textcolor[rgb]{0.53,0.00,0.00}{##1}}}
\@namedef{PY@tok@nl}{\def\PY@tc##1{\textcolor[rgb]{0.46,0.46,0.00}{##1}}}
\@namedef{PY@tok@ni}{\let\PY@bf=\textbf\def\PY@tc##1{\textcolor[rgb]{0.44,0.44,0.44}{##1}}}
\@namedef{PY@tok@na}{\def\PY@tc##1{\textcolor[rgb]{0.41,0.47,0.13}{##1}}}
\@namedef{PY@tok@nt}{\let\PY@bf=\textbf\def\PY@tc##1{\textcolor[rgb]{0.00,0.50,0.00}{##1}}}
\@namedef{PY@tok@nd}{\def\PY@tc##1{\textcolor[rgb]{0.67,0.13,1.00}{##1}}}
\@namedef{PY@tok@s}{\def\PY@tc##1{\textcolor[rgb]{0.73,0.13,0.13}{##1}}}
\@namedef{PY@tok@sd}{\let\PY@it=\textit\def\PY@tc##1{\textcolor[rgb]{0.73,0.13,0.13}{##1}}}
\@namedef{PY@tok@si}{\let\PY@bf=\textbf\def\PY@tc##1{\textcolor[rgb]{0.64,0.35,0.47}{##1}}}
\@namedef{PY@tok@se}{\let\PY@bf=\textbf\def\PY@tc##1{\textcolor[rgb]{0.67,0.36,0.12}{##1}}}
\@namedef{PY@tok@sr}{\def\PY@tc##1{\textcolor[rgb]{0.64,0.35,0.47}{##1}}}
\@namedef{PY@tok@ss}{\def\PY@tc##1{\textcolor[rgb]{0.10,0.09,0.49}{##1}}}
\@namedef{PY@tok@sx}{\def\PY@tc##1{\textcolor[rgb]{0.00,0.50,0.00}{##1}}}
\@namedef{PY@tok@m}{\def\PY@tc##1{\textcolor[rgb]{0.40,0.40,0.40}{##1}}}
\@namedef{PY@tok@gh}{\let\PY@bf=\textbf\def\PY@tc##1{\textcolor[rgb]{0.00,0.00,0.50}{##1}}}
\@namedef{PY@tok@gu}{\let\PY@bf=\textbf\def\PY@tc##1{\textcolor[rgb]{0.50,0.00,0.50}{##1}}}
\@namedef{PY@tok@gd}{\def\PY@tc##1{\textcolor[rgb]{0.63,0.00,0.00}{##1}}}
\@namedef{PY@tok@gi}{\def\PY@tc##1{\textcolor[rgb]{0.00,0.52,0.00}{##1}}}
\@namedef{PY@tok@gr}{\def\PY@tc##1{\textcolor[rgb]{0.89,0.00,0.00}{##1}}}
\@namedef{PY@tok@ge}{\let\PY@it=\textit}
\@namedef{PY@tok@gs}{\let\PY@bf=\textbf}
\@namedef{PY@tok@gp}{\let\PY@bf=\textbf\def\PY@tc##1{\textcolor[rgb]{0.00,0.00,0.50}{##1}}}
\@namedef{PY@tok@go}{\def\PY@tc##1{\textcolor[rgb]{0.44,0.44,0.44}{##1}}}
\@namedef{PY@tok@gt}{\def\PY@tc##1{\textcolor[rgb]{0.00,0.27,0.87}{##1}}}
\@namedef{PY@tok@err}{\def\PY@bc##1{{\setlength{\fboxsep}{\string -\fboxrule}\fcolorbox[rgb]{1.00,0.00,0.00}{1,1,1}{\strut ##1}}}}
\@namedef{PY@tok@kc}{\let\PY@bf=\textbf\def\PY@tc##1{\textcolor[rgb]{0.00,0.50,0.00}{##1}}}
\@namedef{PY@tok@kd}{\let\PY@bf=\textbf\def\PY@tc##1{\textcolor[rgb]{0.00,0.50,0.00}{##1}}}
\@namedef{PY@tok@kn}{\let\PY@bf=\textbf\def\PY@tc##1{\textcolor[rgb]{0.00,0.50,0.00}{##1}}}
\@namedef{PY@tok@kr}{\let\PY@bf=\textbf\def\PY@tc##1{\textcolor[rgb]{0.00,0.50,0.00}{##1}}}
\@namedef{PY@tok@bp}{\def\PY@tc##1{\textcolor[rgb]{0.00,0.50,0.00}{##1}}}
\@namedef{PY@tok@fm}{\def\PY@tc##1{\textcolor[rgb]{0.00,0.00,1.00}{##1}}}
\@namedef{PY@tok@vc}{\def\PY@tc##1{\textcolor[rgb]{0.10,0.09,0.49}{##1}}}
\@namedef{PY@tok@vg}{\def\PY@tc##1{\textcolor[rgb]{0.10,0.09,0.49}{##1}}}
\@namedef{PY@tok@vi}{\def\PY@tc##1{\textcolor[rgb]{0.10,0.09,0.49}{##1}}}
\@namedef{PY@tok@vm}{\def\PY@tc##1{\textcolor[rgb]{0.10,0.09,0.49}{##1}}}
\@namedef{PY@tok@sa}{\def\PY@tc##1{\textcolor[rgb]{0.73,0.13,0.13}{##1}}}
\@namedef{PY@tok@sb}{\def\PY@tc##1{\textcolor[rgb]{0.73,0.13,0.13}{##1}}}
\@namedef{PY@tok@sc}{\def\PY@tc##1{\textcolor[rgb]{0.73,0.13,0.13}{##1}}}
\@namedef{PY@tok@dl}{\def\PY@tc##1{\textcolor[rgb]{0.73,0.13,0.13}{##1}}}
\@namedef{PY@tok@s2}{\def\PY@tc##1{\textcolor[rgb]{0.73,0.13,0.13}{##1}}}
\@namedef{PY@tok@sh}{\def\PY@tc##1{\textcolor[rgb]{0.73,0.13,0.13}{##1}}}
\@namedef{PY@tok@s1}{\def\PY@tc##1{\textcolor[rgb]{0.73,0.13,0.13}{##1}}}
\@namedef{PY@tok@mb}{\def\PY@tc##1{\textcolor[rgb]{0.40,0.40,0.40}{##1}}}
\@namedef{PY@tok@mf}{\def\PY@tc##1{\textcolor[rgb]{0.40,0.40,0.40}{##1}}}
\@namedef{PY@tok@mh}{\def\PY@tc##1{\textcolor[rgb]{0.40,0.40,0.40}{##1}}}
\@namedef{PY@tok@mi}{\def\PY@tc##1{\textcolor[rgb]{0.40,0.40,0.40}{##1}}}
\@namedef{PY@tok@il}{\def\PY@tc##1{\textcolor[rgb]{0.40,0.40,0.40}{##1}}}
\@namedef{PY@tok@mo}{\def\PY@tc##1{\textcolor[rgb]{0.40,0.40,0.40}{##1}}}
\@namedef{PY@tok@ch}{\let\PY@it=\textit\def\PY@tc##1{\textcolor[rgb]{0.24,0.48,0.48}{##1}}}
\@namedef{PY@tok@cm}{\let\PY@it=\textit\def\PY@tc##1{\textcolor[rgb]{0.24,0.48,0.48}{##1}}}
\@namedef{PY@tok@cpf}{\let\PY@it=\textit\def\PY@tc##1{\textcolor[rgb]{0.24,0.48,0.48}{##1}}}
\@namedef{PY@tok@c1}{\let\PY@it=\textit\def\PY@tc##1{\textcolor[rgb]{0.24,0.48,0.48}{##1}}}
\@namedef{PY@tok@cs}{\let\PY@it=\textit\def\PY@tc##1{\textcolor[rgb]{0.24,0.48,0.48}{##1}}}

\def\PYZbs{\char`\\}
\def\PYZus{\char`\_}
\def\PYZob{\char`\{}
\def\PYZcb{\char`\}}
\def\PYZca{\char`\^}
\def\PYZam{\char`\&}
\def\PYZlt{\char`\<}
\def\PYZgt{\char`\>}
\def\PYZsh{\char`\#}
\def\PYZpc{\char`\%}
\def\PYZdl{\char`\$}
\def\PYZhy{\char`\-}
\def\PYZsq{\char`\'}
\def\PYZdq{\char`\"}
\def\PYZti{\char`\~}
% for compatibility with earlier versions
\def\PYZat{@}
\def\PYZlb{[}
\def\PYZrb{]}
\makeatother


    % For linebreaks inside Verbatim environment from package fancyvrb.
    \makeatletter
        \newbox\Wrappedcontinuationbox
        \newbox\Wrappedvisiblespacebox
        \newcommand*\Wrappedvisiblespace {\textcolor{red}{\textvisiblespace}}
        \newcommand*\Wrappedcontinuationsymbol {\textcolor{red}{\llap{\tiny$\m@th\hookrightarrow$}}}
        \newcommand*\Wrappedcontinuationindent {3ex }
        \newcommand*\Wrappedafterbreak {\kern\Wrappedcontinuationindent\copy\Wrappedcontinuationbox}
        % Take advantage of the already applied Pygments mark-up to insert
        % potential linebreaks for TeX processing.
        %        {, <, #, %, $, ' and ": go to next line.
        %        _, }, ^, &, >, - and ~: stay at end of broken line.
        % Use of \textquotesingle for straight quote.
        \newcommand*\Wrappedbreaksatspecials {%
            \def\PYGZus{\discretionary{\char`\_}{\Wrappedafterbreak}{\char`\_}}%
            \def\PYGZob{\discretionary{}{\Wrappedafterbreak\char`\{}{\char`\{}}%
            \def\PYGZcb{\discretionary{\char`\}}{\Wrappedafterbreak}{\char`\}}}%
            \def\PYGZca{\discretionary{\char`\^}{\Wrappedafterbreak}{\char`\^}}%
            \def\PYGZam{\discretionary{\char`\&}{\Wrappedafterbreak}{\char`\&}}%
            \def\PYGZlt{\discretionary{}{\Wrappedafterbreak\char`\<}{\char`\<}}%
            \def\PYGZgt{\discretionary{\char`\>}{\Wrappedafterbreak}{\char`\>}}%
            \def\PYGZsh{\discretionary{}{\Wrappedafterbreak\char`\#}{\char`\#}}%
            \def\PYGZpc{\discretionary{}{\Wrappedafterbreak\char`\%}{\char`\%}}%
            \def\PYGZdl{\discretionary{}{\Wrappedafterbreak\char`\$}{\char`\$}}%
            \def\PYGZhy{\discretionary{\char`\-}{\Wrappedafterbreak}{\char`\-}}%
            \def\PYGZsq{\discretionary{}{\Wrappedafterbreak\textquotesingle}{\textquotesingle}}%
            \def\PYGZdq{\discretionary{}{\Wrappedafterbreak\char`\"}{\char`\"}}%
            \def\PYGZti{\discretionary{\char`\~}{\Wrappedafterbreak}{\char`\~}}%
        }
        % Some characters . , ; ? ! / are not pygmentized.
        % This macro makes them "active" and they will insert potential linebreaks
        \newcommand*\Wrappedbreaksatpunct {%
            \lccode`\~`\.\lowercase{\def~}{\discretionary{\hbox{\char`\.}}{\Wrappedafterbreak}{\hbox{\char`\.}}}%
            \lccode`\~`\,\lowercase{\def~}{\discretionary{\hbox{\char`\,}}{\Wrappedafterbreak}{\hbox{\char`\,}}}%
            \lccode`\~`\;\lowercase{\def~}{\discretionary{\hbox{\char`\;}}{\Wrappedafterbreak}{\hbox{\char`\;}}}%
            \lccode`\~`\:\lowercase{\def~}{\discretionary{\hbox{\char`\:}}{\Wrappedafterbreak}{\hbox{\char`\:}}}%
            \lccode`\~`\?\lowercase{\def~}{\discretionary{\hbox{\char`\?}}{\Wrappedafterbreak}{\hbox{\char`\?}}}%
            \lccode`\~`\!\lowercase{\def~}{\discretionary{\hbox{\char`\!}}{\Wrappedafterbreak}{\hbox{\char`\!}}}%
            \lccode`\~`\/\lowercase{\def~}{\discretionary{\hbox{\char`\/}}{\Wrappedafterbreak}{\hbox{\char`\/}}}%
            \catcode`\.\active
            \catcode`\,\active
            \catcode`\;\active
            \catcode`\:\active
            \catcode`\?\active
            \catcode`\!\active
            \catcode`\/\active
            \lccode`\~`\~
        }
    \makeatother

    \let\OriginalVerbatim=\Verbatim
    \makeatletter
    \renewcommand{\Verbatim}[1][1]{%
        %\parskip\z@skip
        \sbox\Wrappedcontinuationbox {\Wrappedcontinuationsymbol}%
        \sbox\Wrappedvisiblespacebox {\FV@SetupFont\Wrappedvisiblespace}%
        \def\FancyVerbFormatLine ##1{\hsize\linewidth
            \vtop{\raggedright\hyphenpenalty\z@\exhyphenpenalty\z@
                \doublehyphendemerits\z@\finalhyphendemerits\z@
                \strut ##1\strut}%
        }%
        % If the linebreak is at a space, the latter will be displayed as visible
        % space at end of first line, and a continuation symbol starts next line.
        % Stretch/shrink are however usually zero for typewriter font.
        \def\FV@Space {%
            \nobreak\hskip\z@ plus\fontdimen3\font minus\fontdimen4\font
            \discretionary{\copy\Wrappedvisiblespacebox}{\Wrappedafterbreak}
            {\kern\fontdimen2\font}%
        }%

        % Allow breaks at special characters using \PYG... macros.
        \Wrappedbreaksatspecials
        % Breaks at punctuation characters . , ; ? ! and / need catcode=\active
        \OriginalVerbatim[#1,codes*=\Wrappedbreaksatpunct]%
    }
    \makeatother

    % Exact colors from NB
    \definecolor{incolor}{HTML}{303F9F}
    \definecolor{outcolor}{HTML}{D84315}
    \definecolor{cellborder}{HTML}{CFCFCF}
    \definecolor{cellbackground}{HTML}{F7F7F7}

    % prompt
    \makeatletter
    \newcommand{\boxspacing}{\kern\kvtcb@left@rule\kern\kvtcb@boxsep}
    \makeatother
    \newcommand{\prompt}[4]{
        {\ttfamily\llap{{\color{#2}[#3]:\hspace{3pt}#4}}\vspace{-\baselineskip}}
    }
    

    
    % Prevent overflowing lines due to hard-to-break entities
    \sloppy
    % Setup hyperref package
    \hypersetup{
      breaklinks=true,  % so long urls are correctly broken across lines
      colorlinks=true,
      urlcolor=urlcolor,
      linkcolor=linkcolor,
      citecolor=citecolor,
      }
    % Slightly bigger margins than the latex defaults
    
    \geometry{verbose,tmargin=1in,bmargin=1in,lmargin=1in,rmargin=1in}
    
    

\begin{document}
    
    \maketitle
    
    

    
    \section{Mathematical Background of Lattices - Introduction to
Lattice-based
Cryptography}\label{mathematical-background-of-lattices---introduction-to-lattice-based-cryptography}

This is the first post in the series: \textbf{Introduction to
Lattice-based Cryptography}. We focus on lattices as an independent
mathematical concept and its properties. However, we focus on aspects
that would be useful in the study of Lattice-based Cryptography.

Throughout this series, we use 2-dimensions and 3-dimensions because
they are easy to work with. In practice we use sufficiently higher
dimensions.

The only prequisite to start this series is to understand the following:
- \(\mathbb{R}\) represents the set of all real numbers. -
\(\mathbb{Z}\) represents the set of all integers. - \(\mathbb{Q}\)
represents the set of all rational numbers. - \(\mathbb{R}^n\)
represents the set of all n-dimensional real vectors - \(\mathbb{Z}^n\)
represents the set of all n-dimensional integer vectors

    \subsection{What is a Lattice?}\label{what-is-a-lattice}

A Lattice \(L\) is an infinite set of vectors \(\mathbf{v}\) constructed
based two condition: \textbf{discreteness} and \textbf{additivity}.

Formally, an \(n\)-dimensional \emph{lattice} \(L\) is a
\textbf{discrete additive subgroup} of \(\mathbb{R}^n\) generated by a
basis \(B = \{b_1, b_2, ..., b_n\}\) as an \textbf{integer linear
combination}:
\[L = L(B):= \{\sum_{i = 1}^{n}z_ib_i: z_i \in \mathbb{Z}\}\]

where \(z_i\) is an arbitrary integer \(\mathbb{Z}\)

\begin{itemize}
\tightlist
\item
  \textbf{discrete}: This means that every vector point
  \(\mathbf{v} \in L\) has some ``neighborhood'' in which \(\mathbf{v}\)
  is the only lattice point. Loosely speaking, this means that, for
  every point \(\mathbf{v}\) there is ``good space'' around it. By using
  an \emph{integer linear combination} and not just \emph{linear
  combination} we have achieve reasonable discreteness between vector
  points. This is shown below using rationals \(\mathbb{Q}\) as a
  counter example.
\item
  \textbf{additive subgroup}: a lattice \(L\) is an additive subgroup if
  it contains identity element \(0 \in \mathbb{R}^n\)(the zero vector),
  and if for any \(\mathbf{v}, \mathbf{w} \in L\), we have
  \(-\mathbf{v}, -\mathbf{w} \in L\) and
  \(\mathbf{v} + \mathbf{w} \in L\)
\end{itemize}

Below, we will see what a lattice is and what it is not.

    \begin{itemize}
\tightlist
\item
  The singleton set \(\{0\} \in \mathbb{R}^n\) is a lattice \(L\)(for
  any positive integer \(n\)). That is, the zero vector in any dimension
  is a lattice.
\end{itemize}

    \begin{center}
    \adjustimage{max size={0.9\linewidth}{0.9\paperheight}}{fundamentals_files/fundamentals_4_0.png}
    \end{center}
    { \hspace*{\fill} \\}
    
    \begin{center}
    \adjustimage{max size={0.9\linewidth}{0.9\paperheight}}{fundamentals_files/fundamentals_4_1.png}
    \end{center}
    { \hspace*{\fill} \\}
    
    \begin{itemize}
\tightlist
\item
  The integers \(\mathbb{Z} \in \mathbb{R}\) form a 1-dimensional
  lattice \(L\).
\end{itemize}

    \begin{center}
    \adjustimage{max size={0.9\linewidth}{0.9\paperheight}}{fundamentals_files/fundamentals_6_0.png}
    \end{center}
    { \hspace*{\fill} \\}
    
    \begin{center}
    \adjustimage{max size={0.9\linewidth}{0.9\paperheight}}{fundamentals_files/fundamentals_6_1.png}
    \end{center}
    { \hspace*{\fill} \\}
    
    \begin{center}
    \adjustimage{max size={0.9\linewidth}{0.9\paperheight}}{fundamentals_files/fundamentals_6_2.png}
    \end{center}
    { \hspace*{\fill} \\}
    
    \begin{center}
    \adjustimage{max size={0.9\linewidth}{0.9\paperheight}}{fundamentals_files/fundamentals_6_3.png}
    \end{center}
    { \hspace*{\fill} \\}
    
    \begin{center}
    \adjustimage{max size={0.9\linewidth}{0.9\paperheight}}{fundamentals_files/fundamentals_6_4.png}
    \end{center}
    { \hspace*{\fill} \\}
    
    \begin{itemize}
\tightlist
\item
  The integer grid \(\mathbb{Z}^n \in \mathbb{R}^n\) is an n-dimensional
  lattice
\end{itemize}

    \begin{center}
    \adjustimage{max size={0.9\linewidth}{0.9\paperheight}}{fundamentals_files/fundamentals_8_0.png}
    \end{center}
    { \hspace*{\fill} \\}
    
    \begin{center}
    \adjustimage{max size={0.9\linewidth}{0.9\paperheight}}{fundamentals_files/fundamentals_8_1.png}
    \end{center}
    { \hspace*{\fill} \\}
    
    \begin{itemize}
\item
  The set \(\{x ∈ \mathbb{Z}^n : \sum_{i=1}^nx_i ∈ 2\mathbb{Z}\}\) is a
  lattice; it is often called the ``checkerboard'' or ``chessboard''
  lattice, especially in two dimensions. It contains all \(n\)-vectors
  of integers \(x = (x_1, x_2,...,x_n) \in \mathbb{Z}\) such that the
  sum of the components of \(x\), i.e \(\sum_{i = 1}^{x_i}\), is an even
  integer. That is, the sum of elements in the vector is an even
  integer.

  The opposite is not true. That is
  \(\{x ∈ \mathbb{Z}^n : \sum_{i=1}^nx_i ∈ 2\mathbb{Z} + 1\}\) is not a
  lattice. Although it is discrete, it is not an additive subgroup of
  \(\mathbb{R}^n\) mainly because it doesn't contain the identity
  element \(0 \in \mathbb{R}^n\)(the all-zeros vector) as shown in
  non-example diagram below. Recall, that a lattice \(L\) is an additive
  subgroup if it contains identity element \(0 \in \mathbb{R}^n\), and
  if any \(\mathbf{v}, \mathbf{w} \in L\), we have
  \(-\mathbf{v}, -\mathbf{w} \in L\) and
  \(\mathbf{v} + \mathbf{w} \in L\).

  \textbf{Example}: \((0,0),(1,1),(2,4),(−3,5),(−2,−2)\)

  \textbf{Non-example(opposite)}: \((1,0),(2,3),(−1,2)\)
\end{itemize}

    \begin{itemize}
\tightlist
\item
  \emph{\textbf{Example 1}: Just even integers}
\end{itemize}

    \begin{center}
    \adjustimage{max size={0.9\linewidth}{0.9\paperheight}}{fundamentals_files/fundamentals_11_0.png}
    \end{center}
    { \hspace*{\fill} \\}
    
    \begin{center}
    \adjustimage{max size={0.9\linewidth}{0.9\paperheight}}{fundamentals_files/fundamentals_11_1.png}
    \end{center}
    { \hspace*{\fill} \\}
    
    \begin{itemize}
\tightlist
\item
  \emph{\textbf{Example 2}: random 2-tuples and 3-tuples with sum of
  even integer}
\end{itemize}

    \begin{center}
    \adjustimage{max size={0.9\linewidth}{0.9\paperheight}}{fundamentals_files/fundamentals_13_0.png}
    \end{center}
    { \hspace*{\fill} \\}
    
    \begin{center}
    \adjustimage{max size={0.9\linewidth}{0.9\paperheight}}{fundamentals_files/fundamentals_13_1.png}
    \end{center}
    { \hspace*{\fill} \\}
    
    \begin{itemize}
\tightlist
\item
  \emph{\textbf{Non-example}: random 2-tuples and 3-tuples with sum of
  odd integer}
\end{itemize}

    \begin{center}
    \adjustimage{max size={0.9\linewidth}{0.9\paperheight}}{fundamentals_files/fundamentals_15_0.png}
    \end{center}
    { \hspace*{\fill} \\}
    
    \begin{center}
    \adjustimage{max size={0.9\linewidth}{0.9\paperheight}}{fundamentals_files/fundamentals_15_1.png}
    \end{center}
    { \hspace*{\fill} \\}
    
    \begin{itemize}
\item
  The rationals \(\mathbb{Q} \subset \mathbb{R}\) do not form a lattice,
  because although they form a subgroup, it is not discrete: there exist
  rational numbers that are arbitrarily close to zero.

  For two arbitrary rational numbers \(r_1\) and \(r_2\), where
  \(r_1 \lt r_2\) there are inifinitely many rational numbers between
  them therefore making it impossible for either \(r_1\) and \(r_2\) to
  be discrete.

  For example, below is a graph of points in \(\mathbb{Q}\) between 1
  and 2.
\end{itemize}

    \begin{center}
    \adjustimage{max size={0.9\linewidth}{0.9\paperheight}}{fundamentals_files/fundamentals_17_0.png}
    \end{center}
    { \hspace*{\fill} \\}
    
    \begin{center}
    \adjustimage{max size={0.9\linewidth}{0.9\paperheight}}{fundamentals_files/fundamentals_17_1.png}
    \end{center}
    { \hspace*{\fill} \\}
    
    \begin{itemize}
\item
  The odd integers \(2\mathbb{Z} + 1\) do not form a lattice, because
  although they are discrete, they do not form a subgroup of
  \(\mathbb{R}\).

  The odd integers \(2\mathbb{Z} + 1\) do not contain 0
\end{itemize}

\begin{center}\rule{0.5\linewidth}{0.5pt}\end{center}

    \subsection{Basis}\label{basis}

    A basis \(B = \{b_1, b_2, ..., b_n\} \subset \mathbb{R}^n\) of a lattice
\(L\) is a set of linearly independent vector whose integer linear
combinations generate the
lattice:\[L = L(B):= \{\sum_{i = 1}^{n}z_ib_i: z_i \in \mathbb{Z}\}\]

Recall two vectors \(b_1\) and \(b_2\) are said to be linearly
independent if \(w_i * b_1 \neq b_2\) and \(w_j * b_2 \neq b_1\) where
\(w_i, w_j\) is an arbitrary integers.

For example, \(b_1 = [1, 2, 3]\) and \(b_2 = [5, 3, 7]\) are linearly
independent while \(b_3 = [1, 2, 3]\) and \(b_4 = [2, 4, 6]\) are not.
Why? \(w_1 * b_3 = b_4\) and \(w_2 * b_4 = b_3\) where \(w_1\) and
\(w_2\) are \(2\) and \(\dfrac{1}{2}\) respectively. But, we can't find
any \(w_i\) for which \(w_i * b_1 = b_2\) and vice versa.

In general, the vectors \(B = \{b_1, b_2, ..., b_n\}\) are said to be
linearly independent if for any vector \(b_i\) in the set, it cannot be
generated by the linear combination of one, some or all of the other
vectors in the set. That is, \(c_1b_1 + c_2b_2 + ... + c_mb_m \neq b_i\)
where \(m\) is an arbitrary natural number.

We can also represent a basis \(B\) as a matrix. This is an \(n*n\)
matrix where the basis vectors are the ordered columns of the matrix.
This is a non-singular matrix.

With this we can represent a lattice
\(L = B * \mathbb{Z}^n = \{Bz: z \in \mathbb{Z}^n \}\)

For example, given a basis \(B = \{b_1, b_2, b_3\}\) where
\(b_1 = [1, 0, 0]\), \(b_2 = [0, 1, 0]\) and \(b_3 = [0, 0, 1]\)
respectively, we have a matrix
\(B = \begin{bmatrix}1 & 0 & 0 \\[0.3em]0 & 1 & 0 \\[0.3em]0 & 0 & 1
\end{bmatrix}\) and the finite set of integer vectors
\(z = \{v_1, v_2, v_3 \}\) where \(v_1 = [1, 2, 5]\),
\(v_2 = [-5, 7, 9]\) and \(v_3 = [-9, 3, 7]\) respectively, we can
generate a lattice \(L\) with three points \([1, 2, 5]\), \([-5, 7, 9]\)
and \([-9, 3, 7]\).

    This way of representing the basis is instrumental in understanding the
next topic: \textbf{unimodular matrix}

    \subsubsection{Unimodular Matrix}\label{unimodular-matrix}

    Bases \(B_1\), \(B_2\) generate the same lattice \(L\) if and only if
there exists a unimodular \(U \in \mathbb{Z}^{n \times n}\) such that
\(B_1 = B_2U\).

An integer matrix is said to be unimodular if it determinant is
\(\pm1\).

This property of two bases being able to generate the same lattice \(L\)
and \emph{some bases been better than others} is what makes
lattice-based cryptography possible. We will take about \textbf{good and
bad bases} below.

We can efficiently test whether two given matrices \(B_1\), \(B_2\)
generate the same lattice, by checking whether \(B_1^{-1}·B_2\)(i.e dot
product) is unimodular.

\begin{center}\rule{0.5\linewidth}{0.5pt}\end{center}

    \subsection{Norms}\label{norms}

    A \emph{norm} or \emph{vector norm} is a generalization of what we would
normally call the \emph{length} or \emph{magnitude} of a vector.

Recall,
\(||{\mathbf{v}}|| = \sqrt{v_1^2 + v_2^2 + v_3^2 + .... + v_n^2}\), the
\textbf{magnitude} or \textbf{length} or \textbf{size} of a vector
\({\mathbf{v}}\). This is a type of norm called the \textbf{Euclidean
norm} \(||\mathbf{v}||_2\).

Formally, a norm on a vector space \(V\) over \(\mathbb{R}\) or
\(\mathbb{C}\) is a function \(||.||:V \to \mathbb{R}\) that satisfies
the properties for all vectors \(\mathbf{v}\), \(\mathbf{w}\) \(\in\)
\(V\) and all scalars \(\alpha\):

\begin{enumerate}
\def\labelenumi{\arabic{enumi}.}
\item
  \textbf{Non-negativity}: \(||\mathbf{v}|| \geq 0\), and
  \(||\mathbf{v}|| = 0 \iff \mathbf{v} = 0\)
\item
  \textbf{Homogeneity}:
  \(||\alpha \mathbf{v}|| = |\mathbf{\alpha}|||\mathbf{v}||\)
\item
  \textbf{Triangle Inequality}:
  \(||\mathbf{v} + \mathbf{w}|| \leq ||\mathbf{v}|| + ||\mathbf{w}||\)
\end{enumerate}

For example, the Euclidean norm for the vector \(\mathbf{v} = [3, 4]\)
is
\(||\mathbf{v}||_2 = \sqrt{3^2 + 4^2} = \sqrt{9 + 16} = \sqrt{25} = 5\).

The graphical representation is displayed below. The length of the arrow
is \(5\).

This graph would help us understand the next topic:
\textbf{\emph{\(i\)}-th successive minimum}

    \begin{center}
    \adjustimage{max size={0.9\linewidth}{0.9\paperheight}}{fundamentals_files/fundamentals_26_0.png}
    \end{center}
    { \hspace*{\fill} \\}
    
    Other norms include the \emph{Manhattan Norm}, \emph{Maximum Norm} and
\emph{p-norm}. For now, we only care about Euclidean norms.

\begin{center}\rule{0.5\linewidth}{0.5pt}\end{center}

    \subsubsection{Successive Minimum}\label{successive-minimum}

    The Euclidean Norm \(||\mathbf{v}||_2\) of any vector \(\mathbf{v}\) in
a lattice \(L\) can be interpreted as the \textbf{radius} \(r\) of a
circle (or sphere in higher dimensions) where the origin \((0, 0)\) of
the graph is the centre of the circle as shown in the diagram below.

    \begin{center}
    \adjustimage{max size={0.9\linewidth}{0.9\paperheight}}{fundamentals_files/fundamentals_30_0.png}
    \end{center}
    { \hspace*{\fill} \\}
    
    Given a lattice \(L\) and a norm(i.e.~the Euclidean norm), the \(i\)-th
successive minimum \(\lambda_i(L)\) is the smallest radius \(r\) such
that \(L\) contains at least \(i\) linearly independent lattice vectors
of norm at most \(r\) within a ball of radius \(r\) centered at the
origin.

In other words, what is the smallest radius \(r\) such that we can find
\(i\) linearly independent vectors
\(\mathbf{v_1}, \mathbf{v_2}, ..., \mathbf{v_i}\) inside the circle(or
sphere) with radius \(r\)? Recall, that we interpret the norm of a
vector as a radius \(r\) when the centre is at the origin \((0, 0)\).

For example, if we generate the vectors \(\mathbf{v_1} = [4, 2]\),
\(\mathbf{v_2} = [-5, -5]\) and \(\mathbf{v_3} = [2, 11]\) given the
basis points \(\mathbf{b_1} = [2, 1]\) and \(\mathbf{b_2} = [1, 3]\)
creating a lattice \(L\), what is the 1st successive minimum
\(\lambda_1(L)\)? Or, What is \(i\)-th successive minimum when \(i\) is
1?

We start by finding the norm of the vectors(including the basis
vectors):

\begin{itemize}
\tightlist
\item
  \(||\mathbf{b_1}||_2 = \sqrt{2^2 + 1^2} = \sqrt{4 + 1} = \sqrt{5} = 2.24\)
\item
  \(||\mathbf{b_2}||_2 = \sqrt{1^2 + 3^2} = \sqrt{1 + 9} = \sqrt{10} = 3.16\)
\item
  \(||\mathbf{v_1}||_2 = \sqrt{4^2 + 2^2} = \sqrt{16 + 4} = \sqrt{20} = 4.47\)
\item
  \(||\mathbf{v_2}||_2 = \sqrt{-5^2 + -5^2} = \sqrt{25 + 25} = \sqrt{50} = 7.07\)
\item
  \(||\mathbf{v_3}||_2 = \sqrt{2^2 + 11^2} = \sqrt{4 + 121} = \sqrt{125} = 11.18\)
\end{itemize}

The smallest norm is \(2.24\) and if we draw a circle using this norm as
the radius, we find the linearly independent vector
\(\mathbf{b_1} = [2, 1]\) within the circle. Therefore, \(2.24\) is the
the 1st successive minimum.

Notice how it is the closest vector to the origin. It is called the
smallest nonzero vector in the lattice. Take note of this, we discuss
this later in the series.

    \begin{center}
    \adjustimage{max size={0.9\linewidth}{0.9\paperheight}}{fundamentals_files/fundamentals_32_0.png}
    \end{center}
    { \hspace*{\fill} \\}
    
    \subsubsection{Good and Bad Bases}\label{good-and-bad-bases}

Given two bases \(B_1 = \{ [1, 0], [0, 1] \}\) and
\(B_2 = \{[10, 1], [9, 1] \}\) that generate the same lattice(see
\emph{unimodular matrix} above), below they are labelled on the lattice
they generate respectively.

    \begin{center}
    \adjustimage{max size={0.9\linewidth}{0.9\paperheight}}{fundamentals_files/fundamentals_34_0.png}
    \end{center}
    { \hspace*{\fill} \\}
    
    \begin{center}
    \adjustimage{max size={0.9\linewidth}{0.9\paperheight}}{fundamentals_files/fundamentals_34_1.png}
    \end{center}
    { \hspace*{\fill} \\}
    
    If you are given basis \(B_1\) and this lattice, and you are asked to
find the first successive minimum \(\lambda_{1}\), it's very easy to see
that \(B_1\) is the answer because it obviously has the smallest norm.
But then, what if you are given only \(B_2\) and this lattice and you
asked to find the first successive minimum \(\lambda_1\), how hard could
this be?

You can already see that this is a search problem and it gets harder to
solve as the dimension increases. This is in two dimensions so it's easy
to solve but in practice we work in extremely higher dimensions.

That is the point of good and bad bases. Given a good basis it's easy to
find \(\lambda_1\) but it's extremely hard to find given a bad basis.

Formally, we check that a basis is good or bad using the following:

\begin{itemize}
\tightlist
\item
  \textbf{Good bases are close to being orthogonal}: Two vectors
  \(\mathbf{v}, \mathbf{w}\) are orthogonal if their dot product is zero
  (i.e \(\mathbf{v} . \mathbf{w} = 0\)). For example:

  \begin{itemize}
  \tightlist
  \item
    For \(B_1\), \([1, 0].[0, 1] = 0\). This is called \textbf{perfect
    orthogonality}.
  \item
    For \(v_1 = [2, 1]\) and \(v_2 = [1, 3]\), \(v_1.v_2 = 5\). This is
    not the ``best basis'' but it's nothing close to a bad basis. Wait
    until you see that of \(B_2\)
  \item
    For \(B_2\), \([10, 1].[9, 1] = 91\). Now, this is not close to
    being orthogonal.
  \end{itemize}
\item
  \textbf{Bad bases have large Gram-Schmidt Norms}: I won't go into
  details of Gram-Schmidt Norms but keep this in mind.
\end{itemize}

Back to understanding successive minimum.

    \subsubsection{Successive Minimum
(Cont'd)}\label{successive-minimum-contd}

    Again, what is the 2nd successive minimum? We are looking for the
smallest radius(i.e norm) that creates a circle containing two linearly
independent vectors.

From the computed norms, we know that the second smallest vector is
\(\mathbf{b_2} = [1, 3]\) with norm of \(3.16\). Drawing a circle using
this norm as the radius, the circle contains \(\mathbf{b_1}\) and
\(\mathbf{b_2}\). Therefore, \(3.16\) is the 2nd successive minimum.

    \begin{center}
    \adjustimage{max size={0.9\linewidth}{0.9\paperheight}}{fundamentals_files/fundamentals_38_0.png}
    \end{center}
    { \hspace*{\fill} \\}
    
    Thirdly, what is the 3rd successive minimum?

Probably \(4.47\), the third smallest norm. It turns out this is not the
case. They are not linearly independent because \(\mathbf{v_1}\) was
generated by \(b_1\) and \(b_2\) therefore there are only two linearly
independent vectors(i.e \(\mathbf{b_1}\) and \(\mathbf{b_2}\)) in the
circle with radius \(4.47\).

    \begin{center}
    \adjustimage{max size={0.9\linewidth}{0.9\paperheight}}{fundamentals_files/fundamentals_40_0.png}
    \end{center}
    { \hspace*{\fill} \\}
    
    How about \(7.07\)? The fourth smallest norm. This is still not the 3rd
successive minimum for the same reason as before, \(\mathbf{v_2}\) was
generated by \(\mathbf{b_1}\) and \(\mathbf{b_2}\) therefore the circle
with radius \(7.07\) contains only two linearly independent vectors.

    \begin{center}
    \adjustimage{max size={0.9\linewidth}{0.9\paperheight}}{fundamentals_files/fundamentals_42_0.png}
    \end{center}
    { \hspace*{\fill} \\}
    
    And, lastly, how about \(11.18\)? The largest norm. Is it the 3rd
successive minimum? No, it's not. For the same reason as before.

    \begin{center}
    \adjustimage{max size={0.9\linewidth}{0.9\paperheight}}{fundamentals_files/fundamentals_44_0.png}
    \end{center}
    { \hspace*{\fill} \\}
    
    With this, it is safe to say the 3rd successive minimum and above does
not exist because all other vectors were generated by the basis points
\(\mathbf{b_1}\) and \(\mathbf{b_2}\).

In general, the maximum \(i\) in the \(i\)-th successive minimum is the
number of basis vectors.

It's important to point out that a lattice \(L\) might have multiple
basis vectors and this is an instrumental property of lattices that
makes certain types of lattice-based cryptography possible.

\begin{center}\rule{0.5\linewidth}{0.5pt}\end{center}

    \subsection{Fundamental domains and
parallelepipeds}\label{fundamental-domains-and-parallelepipeds}

In order to understand fundamental domain and parallelepipeds, here's a
quick crash course on abstract algebra.

A \textbf{group} is a set \(G\) with an operation \(*\) that satisfies
four properties: - \textbf{Closure}: if \(a\), \(b \in G\), then
\(a * b \in G\) - \textbf{Associativity}: \((a * b) * c = a * (b * c)\)
for all \(a, b, c \in G\) - \textbf{Identity Element}: There exists an
element \(i \in G\) such that \(i * a = a * i= a\) for all \(a \in G\) -
\textbf{Inverse Element}: For every \(a \in G\), there exists an element
\(a^{-1} \in G\) such that \(a * a^{-1} = a^{-1} * a = i\)(\emph{the
identity element})

Example of a group: - The set of integers \(\mathbb{Z}\) under the
operation, addition \(+\) forms a group: - \textbf{Closure:} if
\(a, b \in \mathbb{Z}\), then \(a + b \in \mathbb{Z}\). For example,
\(1 + 5 = 6\) - \textbf{Associativity}: \((a + b) + c = a + (b + c)\).
For example, \((5 + 3) + 9 = 5 + (3 + 9)\) - \textbf{Identity}: The
number \(0\) satisfies \(0 + a = a + 0 = a\). \(0\) is the identity
element. - \textbf{Inverses}: Each \(a \in \mathbb{Z}\) has an inverse
\(-a\), since \(a + (-a) = 0\). For example, \(3 + (-3) = 0\)

A \textbf{subgroup} \(H\) of a \textbf{group} \(G\) is a subset of \(G\)
that is itself a group under the same operation. That is: -
\textbf{Closure}: if \(a\), \(b \in H\), then \(a * b \in H\) -
\textbf{Identity Element}: The identity element of the group \(G\) is in
\(H\) - \textbf{Inverse Element}: If \(a \in H\), then \(a^{-1} \in H\)

Example of a subgroup: - The set of integers divisible by \(5\)
\(5\mathbb{Z} = \{..., -15, -10, -5, 0, 5, 10, 15,...\}\)(i.e multipling
\(5\) by all integers) is a subgroup of \(\mathbb{Z}\) under addition
because: - \textbf{Closure:} The sum of two numbers divisible by \(5\)
is divisible by \(5\) making it closed under addition. -
\textbf{Identity}: \(0\) is in \(5\mathbb{Z}\) - \textbf{Inverses}:
Every \(a \in 5\mathbb{Z}\) has an inverse \(-a\). For example,
\(10 + (-10) = 0\)

\textbf{Note}: \textbf{Associativity} is automatically holds for
subgroups. Ponder on why this is the case.

Let's do a little experiment!

Imagine we added the numbers \(0\) to \(9\) to the elements in
\(5\mathbb{Z}\) to create new sets and displayed them in two halves.
That is:

First half: -
\(\underline{0} = 0 + 5\mathbb{Z} = \lbrace..., -15, -10, -5, 0, 5, 10, 15,...\rbrace\)
-
\(\underline{1} = 1 + 5\mathbb{Z} = \{..., -14, -9, -4, 1, 6, 11, 16,...\}\)
-
\(\underline{2} = 2 + 5\mathbb{Z} = \{..., -13, -8, -3, 2, 7, 12, 17,...\}\)
-
\(\underline{3} = 3 + 5\mathbb{Z} = \{..., -12, -7, -2, 3, 8, 13, 18,...\}\)
-
\(\underline{4} = 4 + 5\mathbb{Z} = \{..., -11, -6, -1, 4, 9, 14, 19,...\}\)

Second half: -
\(\underline{5} = 5 + 5\mathbb{Z} = \{..., -10, -5, 0, 5, 10, 15, 20,...\}\)
-
\(\underline{6} = 6 + 5\mathbb{Z} = \{..., -9, -4, 1, 6, 11, 16, 21,...\}\)
-
\(\underline{7} = 7 + 5\mathbb{Z} = \{..., -8, -3, 2, 7, 12, 17, 22,...\}\)
-
\(\underline{8} = 8 + 5\mathbb{Z} = \{..., -7, -2, 3, 8, 13, 18, 23,...\}\)
-
\(\underline{9} = 9 + 5\mathbb{Z} = \{..., -6, -1, 4, 9, 14, 19, 24,...\}\)

Do you notice anything about these sets!? Here are some observations:

\begin{enumerate}
\def\labelenumi{\arabic{enumi}.}
\tightlist
\item
  The sets in the second half are a repetition of the sets in the first
  half so we will only focus on the first half moving on. Try creating
  more sets with numbers greater than \(9\) or less than \(0\).
\item
  The sets are disjoint(i.e, they have no elements in common). Every
  element appears in only one set.
\item
  The ``combination'' of the sets makes up the set of integers
  \(\mathbb{Z}\). That is,
  \(\underline{0} \cup \underline{1} \cup \underline{2} \cup \underline{3} \cup \underline{4} = \mathbb{Z}\)
  where \(\cup\) represents union.
\item
  Each set contains element with the same remainder when divided by
  \(5\). For example, every element in \(\underline{2}\) would result in
  a remainder of \(2\) when divided by \(5\). These sets are called
  \textbf{equivalent classes} in modulo arithmetic. In this case, they
  are equivalent because they share the same remainders but equivalence
  relations can be of different forms as we will see later.
\item
  If we add any element from a set to any element from another set, you
  get an element from a set. For example, if we add any number in
  \(\underline{0}\) to any number in \(\underline{1}\), we get a number
  in \(\underline{1}\) and if we add any number in \(\underline{2}\) to
  any number in \(\underline{4}\), we get a number in \(\underline{1}\).
\item
  If we add any number from a set to any number in \(\underline{0}\), we
  get back a number in the same set. For example, if we add \(14\) from
  \(\underline{4}\) to \(15\) from \(\underline{0}\), we get \(4\) from
  \(\underline{4}\).
\item
  If we add any number from \(\underline{2}\) to any number from
  \(\underline{3}\), we get a number from \(\underline{0}\), if we add
  any number from \(\underline{1}\) to any number from
  \(\underline{4}\), we get a number from \(\underline{0}\) and if we
  add any number from \(\underline{0}\) to any number from itself, we
  get a number from \(\underline{0}\).
\item
  The last three observations implies \textbf{closure}, presence of an
  \textbf{identity element} and presence of an \textbf{inverse element}
  respectively. The sets are called \textbf{cosets} and they form a
  \textbf{group} called the \textbf{quotient group}
  \(\mathbb{Z}/5\mathbb{Z} = \{\underline{0}, \underline{1}, \underline{2}, \underline{3}, \underline{4}\}\).
  We call it a quotient group because we use a subgroup to divide the
  group into cosets. Not all subgroups create cosets that form a
  quotient group like \(5\mathbb{Z}\), but when they do we call them
  \textbf{normal subgroups}. Our observations doesn't show
  \textbf{associativity} but it holds and I urge you to try it out
\end{enumerate}

Formally, a \textbf{coset} is formed when you take a subgroup \(H\) of a
group \(G\) and partition \(G\) into sets of the form
\[gH = \{g * h | h \in H\}\] where \(g \in G\) and \(*\) is an
operation.

\textbf{Note}: The cosets from \(\underline{1}\) to \(\underline{4}\)
are \textbf{not} subgroups of \(\mathbb{Z}\) under addition. They don't
contain the identity element, they are not closed under addition and
they don't contain their inverses. Using \(\underline{3}\) as an
example: - \textbf{identity element}: \(0\), the identity element is not
in the set. - \textbf{inverse element}: It doesn't contain the identity
element so it cannot have an inverse element. - \textbf{closure}:
\(3 + 8 = 11\) but \(11\) is not in \(\underline{3}\).

Now, back to lattices. How does all these help in understanding
fundamental domain and parallelepipeds?

The above example shows we could take any arbitrary integer
\(\mathbb{Z}\) and we \textbf{must} be able to represent it using one of
the elements of the quotient group \(\mathbb{Z}/5\mathbb{Z}\). A
\textbf{fundamental domain} is synonymous to this in lattices.

Since a lattice \(L\) is a subgroup of \(\mathbb{R}^n\), we can define
the quotient group \(\mathbb{R}^n/L\), which divides \(\mathbb{R}^n\)
into cosets of \(L\) in \(\mathbb{R}^n\). A quotient group
\(\mathbb{R}^n/L\) of this form corresponds to a fundamental domain of
\(L\). A fundamental domain is a region in \(\mathbb{R}^n\) that
contains exactly one representative from each coset just like
\(\mathbb{Z}/5\mathbb{Z}\). Similar to how \(\mathbb{Z}/5\mathbb{Z}\)
encompasses the set of integers \(\mathbb{Z}\) using five cosets, the
fundamental domain encompasses all the elements of \(\mathbb{R}^n\) by
having a representative for each coset.

Formally, a coset of \(L\) in \(\mathbb{R}^n\) is a set of the form
\(m + L = \{m + \lambda | \lambda \in L \}\) where \(m\) is a vector in
\(\mathbb{R}^n\). If \(m_1, m_2 \in \mathbb{R}^n\), then \(m_1 + L\) and
\(m_2 + L\) are either the same set or disjoint. The set of all such
cosets forms the quotient group \(\mathbb{R}^n/L\).

For example, given a lattice
\(L = \{[5, 2, 1], [2, 1, 3], [4, 5, 1], [7, 2, 1], [10, 7, 7]\}\) and
\(m_1 = [1.5, 2, 3.3]\), \(m_2 = [1, 5.4, 1.1]\) and
\(m_3 = [1/2, 2.5, 3]\), we can form the following cosets:

\begin{itemize}
\tightlist
\item
  \(\underline{m_1} = m_1 + L = [1.5, 2, 3.3] + L = \{[6.5, 4,  4.3], [3.5,  3,  6.3], [5.5,  7,  4.3], [8.5,  4,  4.3], [11.5,  9, 10.3]\}\)
\item
  \(\underline{m_2} = m_2 + L = [1, 5.4, 1.1] + L = \{[6,  7.4,  2.1], [3,  6.4,  4.1], [5, 10.4,  2.1], [8, 7.4,  2.1], [11, 12.4,  8.1]\}\)
\item
  \(\underline{m_3} = m_3 + L = [1/2, 2.5, 3] + L = \{[5.5,  4.5,  4], [2.5,  3.5,  6], [4.5,  7.5,  4], [7.5,  4.5,  4], [10.5,  9.5, 10]\}\)
\end{itemize}

An example of a fundamental domain is the set generated by
\(F = \{t_1v_1 + t_2v_2 + ... + t_nv_n\ : 0 \leq t_i \lt 1 \}\) where
\(v_1, v_2, ..., v_n\) is the basis of the a lattice and \(t_i\) is a
real number \(\mathbb{R}\) not an integer \(\mathbb{Z}\). This
fundamental domain is called a \textbf{fundamental parallelepiped}.

We are saying that is for any vector \(s \in \mathbb{R}^n\), we would be
able to find its representative in the fundamental parallelepipeds.

Let's look at an example using basis \(v_1 = [2, 1]\) and
\(v_2 = [1, 3]\).

First, the diagram below shows the lattice generate by this basis.

    \begin{center}
    \adjustimage{max size={0.9\linewidth}{0.9\paperheight}}{fundamentals_files/fundamentals_47_0.png}
    \end{center}
    { \hspace*{\fill} \\}
    
    Next, let's create the fundamental parallelepiped and show it.

As stated above, the fundamental parallelepiped would be generated using
\(F = \{t_1v_1 + t_2v_2 : 0 \leq t_i \lt 1 \}\). There are infinitely
many real numbers in the interval \(0 \leq t_i \lt 1\) so we will show
generate a few and show more in the diagram.

\begin{itemize}
\tightlist
\item
  \(t_1 = 0, t_2 = 0\), then
  \(t_1v_1 + t_2v_2 = 0 * [2, 1] + 0 * [1, 3] = [0, 0]\)
\item
  \(t_1 = 0, t_2 = 0.1\), then
  \(t_1v_1 + t_2v_2 = 0 * [2, 1] + 0.1 * [1, 3] = [0.1, 0.3]\)
\item
  \(t_1 = 0.1, t_2 = 0.1\), then
  \(t_1v_1 + t_2v_2 = 0.1 * [2, 1] + 0.1 * [1, 3] = [0.2, 0.1] + [0.1, 0.3] = [0.3, 0.4]\)
\item
  \(...\)
\item
  \(t_1 = 0.9, t_2 = 0.9\), then
  \(t_1v_1 + t_2v_2 = 0.9 * [2, 1] + 0.9 * [1, 3] = [1.8, 0.9] + [0.9, 2.7] = [2.7, 3.6]\)
\end{itemize}

That is, the fundamental parallelepiped
\(F = \{[0, 0], [0.1, 0.3], [0.3, 0.4], ..., [2.7, 3.6]\}\).

Below is the what it looks like on a graph.

    \begin{center}
    \adjustimage{max size={0.9\linewidth}{0.9\paperheight}}{fundamentals_files/fundamentals_49_0.png}
    \end{center}
    { \hspace*{\fill} \\}
    
    Notice how it looks like a scaled down and compressed version of the
lattice.

Now, let's verify our claim that we can find a representative of any
vector \(m \in \mathbb{R}^n\) in the fundamental domain. We will pick
\(s = [5.2, 6.3]\).

This means that there exists \(t_1\) and \(t_2\) such that
\(t_1 * [2, 1] + t_2 * [1, 3] = [5.2, 6.3]\) and we want to find them.
Recall that to create an element of an fundamental parallelepiped in
\(\mathbb{R}^2\) we compute \(t_1v_1 + t_2v_2 = s\). Here, we know
\(s\), we know \(v_1\) and \(v_2\)(the basis) but we don't know \(t_1\)
and \(t_2\) hence we are finding them.

Expanding the expression, we have that \(2t_1 + t_2 = 5.2\) and
\(t_1 + 3t_2 = 6.3\). Solving this we have that \(t_1 = 1.86\) and
\(t_2 = 1.48\)(i.e using any of the methods for solving simultaneous
equations like gaussian elimination).

Recall that \(t_1\) and \(t_2\) are supposed to be in the interval
\(0 \leq t_i \lt 1\). We would have to wrap \(t_1\) and \(t_2\) around
this interval. In other words, we would have to compute \(t_1 \% 1\) and
\(t_2 \% 1\). Doing that we have that \(t_1 = 0.86\) and \(t_2 = 0.48\).

To get to representative \(r\) of \(s\) in the fundamental
parallelepiped, we compute
\(r = t_1v_1 + t_2v_2 = 0.86 * [2, 1] + 0.48 * [1, 3] = [1.72, 0.86] + [0.48, 1.44] = [2.2, 2.3]\).
Therefore, \(r = [2.2, 2.3]\) is the representative of
\(s = [5.2, 6.3]\) in the fundamental parallelepiped.

We are not done yet.

Do you see any similarity between \(r\) and \(m\)? You would notice they
both end in the same decimals \(.2\) and \(.3\). That means their
difference is an integer vector. That is,
\(s - r = [5.2, 6.3] - [2.2, 2.3] = [3, 4]\).

Using our basis \(v_1 = [2, 1]\) and \(v_2 = [1, 3]\), we can generate
the lattice:
\[L = \{[2, 1], [1, 3], [-2 -11], [3, 4], [-11, -13], [9, 7], [1, 8], [4, 7], [-5, -10], [8, 9]\}\]
Recall that, a coset of \(L\) in \(\mathbb{R}^n\) is a set of the form
\(m + L = \{m + \lambda | \lambda \in L \}\). Setting \(r = m\), we have
that:
\[r + L = [2.2, 2.3] + L = \{[4.2,   3.3],[3.2, 5.3], [0.2,  -8.7], [5.2,   6.3], [-8.8, -10.7], [11.2,   9.3], [3.2,  10.3], [6.2,   9.3], [-2.8,  -7.7], [10.2,  11.3]\}\]
Notice how \([3, 4]\) is part of the lattice \(L\) and \([5.2, 6.3]\) is
part of the coset \(r + L\).

This shows that every vector in this coset has the representative
\([2.2, 2.3]\) in the fundamental parallelepiped and this is their
equivalence relation.

Lastly, note that \textbf{a fundamental domain is not a lattice}.

\begin{center}\rule{0.5\linewidth}{0.5pt}\end{center}

    The goal of this material was to introduce these concepts independently;
appreciating their beauty without drawing parallels to their
applications in lattice-based cryptography; we will do that in the other
posts in this series.

In the next post, we will be looking into \emph{Hard Problems in
Lattice-Based Cryptography}. See you there!


    % Add a bibliography block to the postdoc
    
    
    
\end{document}
